\section{Special Abilities}
Special abilities are additional rules that supplement the general rules. A base with an ability that affects other bases, such as \textit{Medic} or \textit{Leader}, cannot benefit from its own ability but may benefit from the abilities of other bases.

Abilities requiring line of sight, as well as \textit{Mechanic} and \textit{Medic}, cannot be used from inside a transport.

\subsection{Psychic Abomination}

Any enemy Psyker within 25~cm of a base with this ability may only use a psychic power after rolling 5+ on 1D6.

Daemons within 25~cm are also more unstable and suffer a $-1$ modifier to their Instability rolls.

Bases with this ability gain a \textit{Psychic Save (2+)}.

\subsection{Agile}

A Titan or Praetorian with this ability is not limited to a total of 90\textdegree{} of turning during each turn.

\subsection{Anti-Aircraft}

These weapons operate independently and may be activated separately from a base's other weapons.

If the base has a First Fire order, these weapons may perform Overwatch Fire against targets at altitude without suffering the usual To-Hit penalty. They may also perform Overwatch Fire while the base has an Advance order, but suffer the normal $-1$ To-Hit penalty.

Weapons with this ability always have a 360\textdegree{} firing arc and reduce a Flyer's Protection save by 2.

When firing at targets that are not at altitude, these weapons function normally but suffer a $-1$ To-Hit penalty.

\subsection{Antigrav}

A base with this ability can only be pinned in an assault by another base with the \textit{Antigrav} or \textit{Jump Packs} ability.

If the Antigrav remains in the assault, either voluntarily or because it has already moved, the normal pinning rules apply.

Antigrav bases ignore enemy zones of control except those generated by bases of the same or a higher class with the \textit{Antigrav} or \textit{Jump Packs} ability. \textit{(Inc)}

They also ignore terrain modifiers while moving but cannot end their movement on Impassable terrain. While moving over terrain features or enemy troops, they are considered to be at altitude.

Antigrav bases may engage Floating bases at altitude in an assault. If they do so, they are considered to be at altitude for the remainder of the turn.

Bases with the \textit{Antigrav} ability may perform Surprise Attacks.

\subsubsection{Surprise Attack}

A base with the \textit{Antigrav} ability and a First Fire order may perform a Surprise Attack.

When activated during the Combat Phase, the base rises to altitude and remains there for the rest of the turn. It returns to ground level during the End-of-Turn Effects step. From its elevated position, it resolves its shooting normally.

If the base is engaged in an assault against bases that cannot follow it to altitude, rising to altitude is treated as disengaging from the assault.

\subsection{Heavy Antigrav}

Bases with the \textit{Heavy Antigrav} ability follow the rules for Antigrav but cannot perform Surprise Attacks.

\subsection{Artillery}

Weapons with the \textit{Artillery} ability fire in a high arc and may shoot without line of sight.

Obstructing terrain crossed by the attack does not provide protection. If the targets are inside the terrain feature, they receive its cover save normally.

Artillery cannot target bases at altitude or perform Overwatch Fire.

Weapons with this ability may perform Indirect Fire against targets the firing base cannot see. To do so, the artillery base must have a First Fire order.

Indirect Fire may be either observed or unobserved. Observed fire requires an unpinned allied base with the \textit{Forward Observer} ability and line of sight to the artillery's intended target. A Forward Observer redirecting fire cannot shoot during the same turn.

\subsubsection{Observed Indirect Fire}

Observed Indirect Fire suffers a $-1$ To-Hit penalty.

\subsubsection{Unobserved Indirect Fire}

Indirect Fire may be performed without line of sight and without a Forward Observer, but suffers a $-2$ To-Hit penalty.

\subsection{Off-Table Artillery}

Off-Table Artillery represents batteries of extremely long-ranged weapons deployed far from the battlefield, including orbital bombardments and naval artillery.

Off-Table Artillery is purchased normally when creating the army list and may only be used once. Different types are listed in the relevant Codex, and the type used is selected when the attack is called.

To call an Off-Table Artillery attack, an unpinned Forward Observer must have line of sight to the targeted point. The attack's Destruction Points are awarded when it is used.

For the purpose of determining the attack's direction, Off-Table Artillery is always considered to originate along the axis of the Forward Observer.

\subsection{Lightning Attack}

A base with this ability may make a 12~cm Take Position move instead of a 6~cm move when resolving an assault.

\subsection{First Strike (X)}

A weapon with this ability may attack a base in contact with it immediately before an assault is resolved.

Resolve the attack after the assault has been selected for activation but before resolving anything else. It is treated as a shooting attack, but its To-Hit roll never suffers penalties and it ignores Shields.

The target may make a normal saving throw. Damage inflicted by First Strike attacks counts towards the final combat result.

If several bases have this ability, the opponent of the player who activated the assault decides first whether to use First Strike. If bases on both sides have this ability, their First Strike attacks are resolved simultaneously.

\subsection{Psychic Attack}

A Psychic Attack can only be negated by a Psychic Save.

Wounds lost to a Psychic Attack cannot be recovered using the \textit{Regeneration} ability. Cover saves cannot be used against this type of power.

If the target has a hit-location chart, the attack always hits its bridge or head location.

\subsection{Authority}

While a base with this ability is present on the battlefield, its controlling player may reroll one Morale Test per turn.

\subsection{Battery}

Weapons with the \textit{Battery} ability combine their fire into a single powerful attack.

All bases in the detachment produce a single template. The attack's To-Hit value depends on the number of bases remaining in the detachment. The To-Hit value shown on the profile applies to the complete detachment.

For each base missing from the detachment, worsen both the attack's To-Hit value and the AP of its \textit{Damages Buildings} ability by 1.

Range and line of sight may be measured from any base in the detachment that is able to fire.

\subsection{All-Round Armour}

A base with this ability does not suffer the additional $-1$ AP modifier for attacks originating within its rear arc.

\subsection{Bombing (X)}

Flyers and Floating bases at altitude may drop X bombs during their movement. Resolve each attack immediately during the movement, following the normal shooting procedure.

Bombs cannot be dropped while the detachment has a Charge order and cannot affect targets at altitude.

If the weapon uses templates, all its templates must touch one another and are treated as a single template. The templates must be centred on the axis of the bombing base's movement.

If the weapon does not use a template, its attacks are made along the axis of movement and may target a detachment over which the base passes.

\subsection{Energy Shields (X)}

When a weapon with at least AP~$-1$ hits a base with this ability, one Energy Shield is destroyed instead of the base suffering a hit. The base begins with X Energy Shields.

An attack with AP~0 cannot damage the base while at least one shield remains active.

Destroyed shields may be repaired during the End-of-Turn Effects step. Roll 1D6 for each destroyed shield; on a result of 4+, that shield is repaired.

While the base has at least one active shield, it also receives a \textit{Psychic Save (4+)}.

\subsection{Bulldozer}

Bases with this ability ignore movement penalties caused by forests, ruins, craters, trenches, Dragon's Teeth, and barricades.

\subsection{Camouflage}

If a base with this ability occupies terrain that provides a cover save, improve that cover save by 1.

The cover save may be improved by a further 1 if the entire detachment has not yet fired and gives up its ability to fire during the current turn.

Camouflage only functions against enemy bases more than 25~cm away. A base engaged in an assault loses the benefit of Camouflage for the remainder of the turn.

\subsection{Advanced Camouflage}

Advanced Camouflage functions like Camouflage, with the following addition:

If the terrain occupied by the detachment does not provide a cover save, the detachment receives a 5+ cover save. This becomes a 4+ cover save if the entire detachment has not yet fired and gives up its ability to fire during the current turn.

\subsection{Devastating Charge}

If a base with this ability is killed by Overwatch Fire while engaging in an assault, it may continue moving by a distance equal to its Movement characteristic, provided this allows it to engage the firing detachment. It then resolves the assault normally during the Combat Phase.

Bases with this ability always resolve their assault during the Combat Phase, even if they were killed before being activated.

A base killed before resolving its assault is removed and counted as a casualty after resolving the assault, regardless of whether its side won.

\subsection{Charismatic (X)}

All allied detachments with at least one base within X~cm of a base with this ability receive a $+1$ bonus to their Morale value.

During Rally Tests in the End Phase, a Rally Test may also be attempted for a detachment within X~cm of the Charismatic base.

\subsection{Daemon Hunter}

Daemon Hunters are troops specialising in combating the daemons of Chaos.

A base with this ability receives a $+1$ AF bonus during an assault if its target has the \textit{Daemon}, \textit{Daemon Prince}, \textit{Greater Daemon}, or \textit{Daemon Engine} ability.

Their holy litanies weaken the daemons' hold upon reality. Every daemon detachment within 12~cm of a Daemon Hunter suffers a $-2$ modifier to its Daemonic Instability Test.

\subsection{Extended Coherency (X)}

The coherency distance of bases with this ability is X~cm instead of 6~cm.

The engagement rules are also different. A detachment engaging a detachment with Extended Coherency does not need to engage all its bases before engaging another detachment.

\subsection{Urban Combat}

When resolving an assault, opponents of a base with this ability receive no AF bonuses from terrain features.

\subsection{Confusion ($-X$)}

A base with this ability may attempt to change the order assigned to an enemy detachment that has at least one base within 25~cm and has not yet been activated during the Movement Phase.

The affected detachment must make a Morale Test with a modifier of $-X$.

If the test is failed, its order is changed. Roll 1D6 to determine its new order:

\begin{center}
    \begin{tabular}{c l}
        \textbf{Roll} & \textbf{New order} \\
        \hline
        1  & First Fire order \\
        2  & Advance order \\
        3  & Forced March order \\
        4  & Charge order \\
        5+ & Chosen by the player using this ability
    \end{tabular}
\end{center}

The new order is shown to both players and then placed face down. The detachment may subsequently be activated normally but cannot perform Overwatch Fire during that turn.

\subsection{Autonomous Defence}

A base with this ability is always considered to have a First Fire order, including immediately after arriving from off-table.

\subsection{Close Defences (X+)}

Every Class 1 or 2 base that engages, or is engaged by, a base with Close Defences suffers a hit on X+ with AP~0.

Resolve the attack when the bases make contact. Cover saves may be used against Close Defences.

\subsection{Damage (+X)}

A weapon with this ability inflicts X additional hits against a base successfully hit by the weapon.

\subsection{Damage (+X) in Assault}

A weapon with this ability inflicts X additional hits after its base wins a duel, in addition to the normal hit.

\subsection{Free Deployment}

Detachments with this ability may deploy anywhere within their controlling player's half of the battlefield.

Free Deployment is resolved at the same time as Infiltration movement. Every base deployed in this manner must be placed at least 12~cm from every enemy base.

Once deployed, the detachment may also use its \textit{Infiltration} ability, if it possesses it.

\subsection{Disruption}

Disruption weapons are powerful enough to penetrate several shields.

When an attack from a Disruption weapon hits a shielded target, it inflicts one hit normally. It then inflicts another hit with its AP and Damage values reduced by 1. Continue inflicting hits in this manner until its AP reaches 0, at which point the attack has no further effect.

If all shields are destroyed before this occurs, resolve one final hit against the target using the remaining AP and Damage values.

\subsection{Front Double Save}

A base with this ability may reroll failed saving throws against attacks originating within its front arc.

\subsection{Duellist}

A base with this ability chooses the order in which its assault duels are resolved.

\subsection{Hard to Hit}

All ranged attacks targeting a base with this ability suffer a $-1$ To-Hit penalty.

This penalty does not apply to template weapons.

\subsection{Dread ($-X$)}

When activated during the Combat Phase, a base with this ability may force one enemy detachment with at least one base within 20~cm to make a Morale Test.

The test suffers a modifier of $-X$. If it is failed, the detachment receives a Fall Back order and immediately makes a Fall Back move.

A detachment cannot be targeted by this ability more than once during the same turn.

\subsection{Elite (X)}

At the beginning of the battle, your army receives a shared pool of Elite rerolls.

Each detachment containing bases with \textit{Elite (X)} adds X rerolls to this pool, regardless of how many Elite bases the detachment contains.

During the battle, Elite rerolls may be spent to reroll dice rolled by your Elite bases. Before using them, declare the total number of dice that will be rerolled.

Only one Elite reroll may be used per base for each dice roll. Consequently, only one die from an assault roll may be rerolled.

Used Elite rerolls are removed from the army's pool.

Elite rerolls may be used for:

\begin{itemize}
    \item To-Hit rolls for any type of shooting attack.
    \item Armour saving throws.
    \item Assault rolls.
    \item Dodge rolls.
    \item Opportunity attacks made when an enemy disengages.
\end{itemize}

\subsection{Damages Buildings (AP \boldmath{$-X$} / Y)}

A base with a weapon possessing this ability may attack destructible terrain features.

When using a template weapon, the centre of the template must be positioned over the structure for it to be hit.

The structure must make Y saving throws with an AP modifier of $-X$. It loses one Wound for each failed save. Saving throws made by structures use the total result of 2D6.

If a detachment has both weapons capable of damaging buildings and ordinary weapons, the ordinary attacks are always resolved first.

\subsection{Damages Buildings (AP \boldmath{$-X$} / Y) in Assault}

A base with a weapon possessing this ability may attack a destructible terrain feature with which it is in contact during the Combat Phase.

If the base is also engaged in an assault against enemy bases, it must choose whether to use the weapon against its opponents or the terrain feature.

The building-damaging action is resolved when the base is activated during the Combat Phase. If the base is engaged in an assault, the attack is resolved after the assault itself, provided the base survives.

Casualties caused by the destruction of the building count towards the assault's combat result.

In all other respects, this ability functions in the same manner as \textit{Damages Buildings}. It cannot be used after a Redeployment action.

\subsection{Entangle}

When a base with this ability makes contact with an enemy base that has a hit-location chart, it may immobilise one of the enemy base's weapons. The effects of that weapon are cancelled for the remainder of the turn.

If two opposing bases both possess such weapons, the abilities cancel one another and both weapons become entangled.

\subsection{Dodge (X+)}

When a base with this ability loses Wounds in an assault, roll one die for each Wound lost. For each result of X+, the base does not lose that Wound.

Dodge may also be used against hits caused by Close Defences.

\subsection{Assault Expert}

A base with this ability wins an assault duel if the final assault roll is tied.

\subsection{Floater}

Floaters follow the rules for Heavy Antigrav.

At the beginning of each turn, a Floater may choose to remain at ground level or rise to altitude. It remains at the chosen level for the entire turn.

While at altitude, it follows the additional rules below.

\subsubsection{Assault}

Only bases with the \textit{Antigrav}, \textit{Heavy Antigrav}, \textit{Floater}, \textit{Flyer}, or \textit{Jump Packs} ability may engage a Floater at altitude in an assault. Any base doing so is considered to be at altitude for the remainder of the turn.

A Floater may engage other troops without restriction. It moves at altitude until directly above its target and then engages it, moving to the target's altitude if necessary.

Floaters may also engage Flyers in an assault.

If a Floater has the \textit{Interceptor} ability and is in a Countercharge position, it may intercept troops moving at altitude within a radius equal to its own Movement characteristic. It then performs an Aerial Interception as described in the Flyer rules.

\subsubsection{Transport}

If a Floater is at altitude, only bases with the \textit{Jump Packs} or \textit{Antigrav} ability may disembark from it. Bases cannot embark while the Floater is at altitude.

\subsubsection{Bombing}

Floaters may use their bombs during movement, but only while at altitude.

\subsection{Deep Strike (X)}

The controlling player selects a point on the battlefield and places one base from the detachment at that point. The base then scatters X times, moving 3D6~cm for each scatter.

If the final point is outside the battlefield, within Impassable terrain, or within the zone of control of an enemy base of the same or a higher class, the detachment does not arrive. Another attempt may be made during the following turn.

Otherwise, place the first base as close as possible to the final arrival point. Place every other base in the detachment anywhere within 6~cm of the first base.

No base may arrive within Impassable terrain or within the zone of control of an enemy base of the same or a higher class.

A detachment entering the battlefield in this manner cannot receive a First Fire order during the turn in which it arrives. It also loses 5~cm from its total available movement during that turn.

\subsubsection{Limited Assault}

On the first turn, a Class 3 or higher detachment cannot select an initial arrival point within the opposing player's half of the battlefield.

\subsection{Advanced Fortification}

Some armies may purchase fortifications. These are deployed before normal deployment and follow the restrictions of \textit{Free Deployment}.

Fortifications cannot be placed within Impassable terrain or on a terrain feature that provides a cover bonus.

If a fortification is attached to a detachment, that detachment must be deployed inside it at this stage. It must receive a First Fire order during the first turn. This order cannot be changed by an allied ability such as \textit{Master Strategist}.

\subsection{Template (X)}

Templates have an area of effect and may therefore affect more than one target.

Any base whose centre is covered by a template may be hit, depending on the template weapon's To-Hit roll.

Obstructing terrain crossed by a template attack does not provide protection. However, a target inside such terrain receives its cover save normally.

A template weapon may target bases at altitude if it fulfils the normal requirements. If it does so, it affects only targets at altitude and cannot affect targets at ground level.

\subsubsection{Structures}

For a circular template to affect a structure and its garrisoned bases, the centre of the template must be positioned within the structure.

For a flame template, the template must cover the centre of the structure.

If these conditions are not met, neither the structure nor its garrisoned bases are affected. Consequently, a single circular template cannot affect more than one structure.

Flame templates and 12~cm circular templates affect every base garrisoned within the structure.

A 7.5~cm circular template affects 50\% of the garrisoned bases, rounded up. If several detachments are garrisoned within the structure, calculate this percentage separately for each detachment.

\subsubsection{Templates That Must Touch}

Some weapons place several templates that must touch one another.

Each additional template must be placed with its edge touching another template created by the same weapon and must satisfy the normal template placement restrictions.

All the templates are then treated as a single large template. Consequently, they all scatter together in the same direction and by the same distance.

If the centre of at least one template covers a structure, its garrisoned bases are affected, but they are treated as being covered by only one template.

\subsubsection{Scatter}

A template that scatters uses a scatter die to determine the direction.

If a ``Hit'' result is rolled, the template does not scatter. Otherwise, it scatters in the direction indicated by the arrow.

When an attack uses several touching templates, place all the templates before making a single scatter roll. Every template then scatters in the same manner.

Flame templates stop if they encounter Occulting terrain of the same class as the firing base, or of a higher class. Bases inside that terrain are still hit, but the template does not extend beyond it.

\subsection{Tactical Genius}

Once per turn, when an allied detachment within 20~cm of a base with this ability is activated, you may change that detachment's order.

\subsection{Graviton}

A Graviton weapon does not have a fixed To-Hit value. Instead, it uses the target's armour Save value as its To-Hit value, to a maximum of 2+.

A base with a 2+ armour Save is therefore hit on a roll of 2+.

A target with a hit-location chart, or a terrain feature, is always hit on a roll of 2+.

Graviton weapons have no effect against a target with at least one active \textit{Energy Shield}, \textit{Energy Field}, \textit{Organic Metal}, or \textit{Deflector Shield}.

\subsection{Ignores Shields}

A weapon with this ability ignores Energy Shields, Energy Fields, Deflector Shields, Organic Metal, and any Psychic Save provided by them.

\subsection{Infiltration}

These bases are stealthy and capable of approaching the enemy before the battle begins.

After deployment, bases with this ability may move up to 25~cm. This movement cannot bring a base into an enemy zone of control and cannot be performed while the detachment is being transported.

\subsection{Engineer}

Engineers are troops specialising in combat engineering.

When Engineers are purchased, they must be added to a Class 1 detachment. That detachment may sacrifice its shooting to perform one of the actions below, provided that it meets all the following requirements:

\begin{itemize}
    \item It is not within an active enemy zone of control.
    \item It is not pinned in an assault and has not engaged in an assault.
    \item It does not have a Charge order and has not performed a Redeployment action.
\end{itemize}

The selected action is resolved when the detachment is activated during the Combat Phase.

\subsubsection{Mines}

The detachment may place a minefield in contact with it. The minefield is 5~cm wide and 15~cm long.

It may be divided into several consecutive sections, but these sections must touch one another and at least one base from the detachment.

Mines cannot be placed beneath a base, within an enemy zone of control, or more than 10~cm from a base in the detachment.

\subsubsection{Building Demolition}

The detachment may inflict one hit against a building with the ability:

\begin{center}
    \textit{Damages Buildings (AP X, 1D3+1)}
\end{center}

X is equal to half the number of bases in the detachment that are in contact with the building, rounded as normal.

\subsubsection{Dig In}

The detachment may dig trenches.

The entire detachment is treated as having a 4+ cover save and a $+2$ AF bonus. These bonuses do not stack with an existing terrain cover bonus and are not lost when making a Consolidation move.

If the detachment moves, it loses the bonuses provided by the trenches.

For the purpose of controlling objectives, the detachment is considered to occupy a terrain feature.

\subsection{Inorganic}

A base with this ability is immune to certain effects. The descriptions of those effects specify when this immunity applies.

\subsection{Interceptor}

A detachment with this ability may perform Aerial Interceptions. See the Flyer rules for further details.

\subsection{Lance (X+ / Y)}

When a base with the \textit{Lance} ability engages a target in an assault, it may make a shooting attack against that target.

Resolve this attack when the attacking base makes contact with its target during the Movement Phase. It may be used during a Charge or Consolidation move.

The target is hit on X+ with AP~Y. The attack otherwise follows the normal shooting rules, including the effects of the \textit{HQ} ability.

Lances ignore Holo-fields, Energy Shields, Energy Fields, and Deflector Shields.

This attack is resolved simultaneously with any Close Defences possessed by the target, even if the target is killed by the Lance.

If the target survives, the base with the Lance stops moving and engages it in an assault.

If the target is killed and has a lower Blocking Class than the base with the Lance, the attacking base may continue moving and engage another base. It cannot use its Lance again during that movement.

\subsection{Robotics Master}

Detachments with the \textit{Robot} ability may activate normally while within 25~cm of a base with this ability.

\subsection{Master Strategist}

If your army contains a base with the \textit{Master Strategist} ability, you may place two orders face down in reserve at the beginning of the battle instead of assigning them to detachments.

During the battle, these reserve orders may be used while at least one base with the Master Strategist ability remains alive.

When activating a Class 1--4 detachment, you may replace its current order with one of the reserve orders. Discard the detachment's original order and assign the reserve order in its place.

If an order is replaced with a First Fire order, the detachment may subsequently perform Overwatch Fire.

Once a reserve order has been assigned, it is removed from the reserve and cannot be used again.

A reserve order cannot be assigned to a detachment that has no order or has a Fall Back order. It also cannot replace an order imposed by another rule or special ability.

\subsection{Walker}

A base with this ability suffers the same terrain movement penalties as a Walker-class base.

\subsection{Mechanic}

Any Walker or Vehicle detachment with at least one base within 6~cm of a base with this ability gains \textit{Regeneration (5+)}.

If the affected base already has Regeneration, improve its Regeneration value by 1. For example, Regeneration (5+) becomes Regeneration (4+).

This ability does not function while the Mechanic is inside a transport. If the Mechanic is garrisoned, the ability can only affect bases garrisoned within the same structure.

Mechanics that are Attached Characters may be attached to a Vehicle detachment. When they are, they are considered to be transported by one of the detachment's bases for the purpose of movement.

\subsection{Medic}

Any Infantry or Cavalry detachment with at least one base within 6~cm of a base with this ability gains \textit{Regeneration (5+)}.

If the affected base already has Regeneration, improve its Regeneration value by 1. For example, Regeneration (5+) becomes Regeneration (4+).

This ability does not function while the Medic is inside a transport. If the Medic is garrisoned, the ability can only affect bases garrisoned within the same structure.

\subsection{Leader}

All allied detachments with at least one base within 12~cm of a base with this ability receive a $+1$ bonus to their AF.

\subsection{Assault Module}

Assault Modules are transports designed specifically for attacking structures.

If a base with this ability is in contact with a structure and troops disembarking from it perform a Charge, those troops receive a $+1$D6 bonus to their AF against bases inside the structure for the remainder of the turn.

In all other respects, an Assault Module functions as a normal transport.

\subsection{Forward Observer (FO)}

A Forward Observer may observe and direct Indirect Artillery Fire.

Forward Observers are also the only bases capable of calling in Off-Table Artillery attacks.

\subsection{Parry}

A base with this ability ignores the first outnumbering die during an assault.

Consequently, the opposing side only begins receiving outnumbering dice when the base fights its third opponent.

\subsection{Character}

Characters represent important individuals who distinguish themselves within an army. This ability is relevant in certain scenarios.

\subsection{Attached Character}

At the beginning of the battle, an Attached Character must join a detachment. The character and the detachment are then treated as a single detachment.

The Attached Character must remain in coherency with the detachment, uses its Morale value, and adds its points cost to the detachment's total cost.

Attached Characters do not occupy space in transports. A single detachment may include up to two Attached Characters.

Whenever possible, an Attached Character purchased as part of a company formation must be attached to a detachment belonging to that company.

An Attached Character gains the Movement characteristic of the detachment to which it is attached. For example, a character attached to a motorcycle detachment gains the Movement characteristic of those motorcycles.

It also gains the detachment's movement-related special abilities, including:

\begin{itemize}
    \item Infiltration
    \item Free Deployment
    \item Jump Packs
    \item Antigrav
    \item Hard to Hit \textit{(Inc)}
    \item Camouflage \textit{(Inc)}
    \item Advanced Camouflage
\end{itemize}

An Infantry Attached Character, Class 1, may be attached to an Infantry, Walker, or Cavalry detachment. If attached to a Walker or Cavalry detachment, it gains that detachment's class.

In all other cases, an Attached Character must join a detachment of the same type.

\subsection{Fear}

A Class 1--3 detachment that enters base-to-base contact with a base possessing this ability must pass a Morale Test or suffer a $-2$ AF penalty for the remainder of the turn.

This applies whether the Fear-causing base charges or is itself charged.

Fear has no effect against bases with the \textit{Fear} or \textit{Terror} ability.

\subsection{Wounds (X)}

A base with this ability has X Wounds, allowing it to survive multiple injuries.

By default, a base has only one Wound.

\subsection{Firing Position (X)}

A base with \textit{Firing Position (X)} \textit{(Inc)} allows transported Class 1 bases to fire while embarked.

The transported bases are activated according to their own orders, independently of the base with the Firing Position ability, and have a 360\textdegree{} firing arc.

Embarked bases using a Firing Position cannot be targeted.

They may fire at enemy bases engaged in an assault with their transport but suffer a $-1$ To-Hit penalty.

Bases using a Firing Position may use their abilities normally, including abilities requiring line of sight, except for the \textit{Medic} and \textit{Mechanic} abilities.

\subsection{Dark Presence (\boldmath{$-X$} / Y cm)}

Enemy detachments with at least one base within Y~cm of a base with this ability suffer a $-X$ modifier to their Morale value.

This effect is not cumulative.

\subsection{Protection (X+)}

Protection is a fixed saving throw made before all other saving throws. It may be used in addition to other types of saving throw.

A base cannot benefit from more than one Protection save. If several are available, use the best one.

\subsection{Psyker}

These troops possess special abilities such as sorcery, mutant powers, or technomancy. Details of their powers are provided in the relevant Codex.

During the Combat Phase, a Psyker may use both one psychic power and its conventional weapons.

Powers marked ``Shooting'' in their description follow the same rules and restrictions as standard shooting attacks.

A power used during the Movement Phase may be used at any point during the Psyker's movement activation, even if the Psyker is pinned in an assault.

Other powers may be used regardless of the Psyker's order and even while the Psyker is engaged in an assault.

Unless otherwise stated, psychic powers have a 360\textdegree{} firing arc.

\subsection{HQ}

HQ bases represent a small number of important individuals. While they remain close to another base of the same Blocking Class, they receive protection against shooting attacks.

When an HQ base is targeted by a shooting attack, its controlling player may select another base of the same class within 6~cm as the attack's target instead.

This decision must be made before any saving throws are rolled.

HQ protection does not function during an assault.

\subsection{Jump Packs}

Bases equipped with Jump Packs may make short flights over terrain, buildings, and enemy troops.

They ignore enemy zones of control except those generated by bases with the \textit{Floater}, \textit{Antigrav}, or \textit{Jump Packs} ability. \textit{(Inc)}

They also ignore terrain modifiers during their movement but cannot end their movement on Impassable terrain.

While using this ability to move over an obstacle, the base is considered to be at altitude.

\subsection{Redeployment (X)}

A base with this ability allows its controlling player to redeploy X detachments at the beginning of the battle.

After both players have finished deploying their armies, but before resolving Infiltration, the player may reposition X detachments according to the normal deployment rules.

\subsection{Reduces Cover (X)}

A weapon with this ability worsens the target's cover save by X.

\subsection{Regeneration (X+)}

When a base with \textit{Regeneration (X+)} would lose one or more Wounds, roll one die for each Wound lost.

For each result equal to or greater than X, the base does not lose that Wound. Regeneration functions against both shooting attacks and assaults.

Regeneration is more difficult during an assault and suffers a $-1$ modifier.

Only one Regeneration attempt may be made for each Wound lost.

Regeneration is resolved before hit-location rolls. Extra damage caused by hit-location effects cannot be regenerated. A successful Regeneration roll prevents the loss of the Wound but does not cancel any additional effects caused by the attack.

\subsection{Reroll 1s in Assault}

A base with this ability may reroll results of 1 on its assault dice. If both sides in an assault can reroll assault dice, the player who did not initiate the assault's resolution rerolls first.

\subsection{Reroll 1s and 2s in Assault}

A base with this ability may reroll results of 1 and 2 on its assault dice. If both sides in an assault can reroll assault dice, the player who did not initiate the assault's resolution rerolls first.

\subsection{Reroll Assault Dice}

A base with this ability may reroll its assault dice. If it does so, it must reroll all of them. If both sides in an assault can reroll assault dice, the player who did not initiate the assault's resolution rerolls first.

\subsection{Robot}

Detachments containing bases with the \textit{Robot} ability receive orders normally, but their activation is handled differently.

If at least one Robot base in the detachment is within 6~cm of an allied non-Robot base that does not have a Fall Back order, the detachment may activate normally.

Otherwise, the detachment must follow these restrictions:

\begin{enumerate}
    \item A Robot detachment with a First Fire order must perform Overwatch Fire whenever it is able to do so.

    \item If a Robot detachment with a First Fire order did not perform Overwatch Fire, it must shoot at the closest visible enemy detachment.

    \item A Robot detachment with an Advance or Forced March order must use its full movement to move towards either the closest enemy detachment or the closest objective, as chosen by the controlling player. It cannot engage in an assault. The Robots may subsequently shoot at targets chosen by the controlling player.

    \item A Robot detachment with a Charge order must move towards the closest enemy detachment and engage it in an assault if possible.
\end{enumerate}

\subsection{Roller (X+ / AP \boldmath{$-Y$})}

A base with this ability and a Charge order may crush bases of a lower class while moving.

When contact is made, the target suffers a hit on X+ with AP~$-Y$. This roll may be affected by the same modifiers as a shooting attack.

If the target survives, the moving base stops and engages it in an assault normally.

If the target is destroyed, the base may continue moving. However, it loses momentum, and its next crushing attack suffers a cumulative $-1$ To-Hit penalty.

The base may continue crushing targets in this manner until the required To-Hit result reaches 7+, at which point it can no longer crush additional bases.

If some bases in a detachment successfully hit with their Roller attacks while others fail, the successful bases may attempt to crush the targets of the bases that failed.

If these targets are destroyed, the bases that initially failed are not considered engaged in an assault and may continue making crushing attacks, including the cumulative penalty for lost momentum.

\subsection{Fearless}

A detachment with this ability automatically passes all Morale Tests.

\subsection{Fixed Save}

A base with a Fixed Save may use it instead of its normal armour save. A Fixed Save is never modified by Armour Penetration.

A Fixed Save is written as X+f in the \textit{Save} column.

If a base has both a normal armour save and a Fixed Save, they are written as X+/Y+f, where X+ is the normal save and Y+f is the Fixed Save.

\subsection{Psychic Save (X+)}

A base with \textit{Psychic Save (X+)} may use this saving throw against psychic powers.

\subsection{Sniper}

When a base with the \textit{Sniper} ability targets an HQ base, roll 1D6.

On a result of 4+, the target cannot use its HQ ability against the Sniper's attack.

\subsection{Terror}

A detachment engaged in an assault by a base with the \textit{Terror} ability must make a Morale Test with a $-1$ modifier.

If the test is failed, the detachment immediately receives a Fall Back order but does not make a Fall Back move. Consequently, it cannot perform Overwatch Fire, even if it would otherwise be able to do so.

A detachment attempting to engage a base with Terror must also pass a Morale Test with a $-1$ modifier.

If this test is failed, the attacking bases stop before entering the Terror-causing base's zone of control. They cannot move any farther but do not receive a Fall Back order.

Bases with the Terror ability are immune to Terror.

\subsection{Fire on the Move}

A base with this ability and an Advance order may move normally.

At the end of its movement, replace its Advance order with an unrevealed First Fire order. It may subsequently perform Overwatch Fire, but doing so imposes an additional $-1$ To-Hit penalty.

Regardless of whether it performs Overwatch Fire, it does not reroll results of 1 when shooting.

\subsection{Reflex Fire}

A base with this ability does not suffer the normal $-1$ To-Hit penalty when performing Overwatch Fire.

\subsection{Subterranean Fire}

Attacks with this ability ignore Energy Shields, Energy Fields, and Deflector Shields.

They have no effect against Antigrav bases or bases at altitude.

Against a base with a hit-location chart, Subterranean Fire always hits its lowest location, Location~1.

\subsection{Open-Topped}

An Open-Topped vehicle has an exposed firing area that allows embarked troops to shoot from the transport.

Embarked bases may shoot at enemy bases engaged in an assault with the Open-Topped transport but suffer a $-1$ To-Hit penalty.

Bases inside an Open-Topped transport are treated as garrisoned and receive a 4+ cover save, but no AF bonus.

Ranged attacks may be allocated freely between the transport and the embarked bases. When a template hits the transport, the embarked troops are also hit.

Embarked bases cannot fight in an assault while inside an Open-Topped transport.

Bases inside an Open-Topped transport may use abilities requiring line of sight, except for the \textit{Medic} and \textit{Mechanic} abilities.

\subsection{Turret}

A weapon with the \textit{Turret} ability has a 360\textdegree{} firing arc.

\subsection{Multiple Hits (X)}

A weapon with this ability has X attack dice and therefore makes X To-Hit rolls against its target. All of the weapon's shots must be directed at the same target.

\subsection{Transport (X)}

A base with \textit{Transport (X)} may transport X Infantry bases.

Entering or leaving a transport costs the transported bases 5~cm of movement. When a base leaves a transport, place it in contact with the transporting base.

If the transport is in contact with a garrison, a transported base may leave the transport and enter the garrison by spending 10~cm of movement.

A base may enter and leave a transport during the same turn. This costs 5~cm to embark and another 5~cm to disembark.

A transport with a capacity of 6 or more may transport Walker bases. Each Walker occupies the capacity of two Infantry bases.

\subsubsection{Encirclement}

Transported bases may disembark even if their transport is engaged in an assault. While disembarking, they ignore the enemy bases and zones of control involved in that assault.

A transport is \textit{Encircled} if the number of enemy bases engaged with it is equal to or greater than its class. Disembarking from an Encircled transport is dangerous.

For each disembarking base, the opposing player may select one base engaged with the transport to make an opportunity attack, following the rules for disengaging from an assault.

Each enemy base may make only one opportunity attack when a detachment disembarks.

\subsubsection{Activations}

During the Movement Phase, a transport and the bases it transports count as a single activation, even if they have different orders.

If a detachment enters a transport, the transport is activated at the same time as the transported detachment.

Conversely, if a transport moves towards a detachment, that detachment may be activated immediately to embark. The transport may then continue moving and may subsequently allow the detachment to disembark.

If either the transport or the transported troops have a First Fire order, that order may remain unrevealed while the other group is activated normally.

If a transport carries several detachments, only the transport and one transported detachment are activated together. The other detachments must be activated later.

Several detachments may disembark during a transport's movement, but only one of them may activate at that time. Place the other detachments along the transport's route with their orders face down. They must be activated later and are considered to have disembarked from a transport.

\subsubsection{Removing Casualties}

When a transport detachment suffers casualties, the player suffering the casualties chooses which transport is destroyed.

This is particularly important if the transports carry different types of bases.

\subsubsection{Abandoning Troops}

If there is insufficient transport capacity for an entire detachment, part of the detachment may be abandoned. The abandoned bases are immediately removed as casualties. Transports may also be abandoned.

\subsubsection{Morale}

While a detachment is embarked within a transport detachment, Morale Tests are made as though the transports and transported bases formed a single detachment.

Use the better Morale value of the transport detachment and the transported detachment.

If the Morale Test is failed, both detachments receive a Fall Back order.

Bases cannot embark or disembark during a Fall Back move.

\subsubsection{Emergency Exit}

Bases inside a destroyed transport must make an Emergency Exit Test.

Roll 1D6 for each base inside the transport. On a result of 2+, modified by the AP of the weapon that destroyed the transport, the base escapes in time. If several AP values apply, use the best one.

If the test is failed, the base is destroyed as though it had been hit by a shooting attack. Abilities such as \textit{HQ} and \textit{Regeneration} may therefore be used.

After removing the destroyed transport, place the surviving bases within the transport's former zone of control. If there is insufficient room, any bases that cannot be placed are removed as casualties.

If the transport was destroyed in an assault or by an attack that does not permit a saving throw, the Emergency Exit Test is an unmodifiable 4+.

\subsection{Attached Transport}

Attached Transports and the bases they transport are treated as a single detachment.

The transports use the Morale value of the detachment to which they are attached.

The transport group and transported group receive separate orders but are activated simultaneously. They have \textit{Extended Coherency (25~cm)} with one another.

The transported troops may begin the battle either embarked or outside their transports.

If Attached Transports carry more than one detachment, they are attached to only one of those detachments. All other transported detachments follow the normal Transport rules.

\subsection{Strategic Vision (X)}

Once per battle, the controlling player may modify their Initiative roll by up to X.

This modifier is applied after both players have rolled for Initiative.

\subsection{Flyer}

Flyers operate differently from other types of unit.

\subsubsection{Deployment}

Flyers may either deploy normally on the battlefield or begin the battle off-table.

A Flyer placed off-table cannot enter the battlefield before Turn 2. When it enters, it begins its movement from the battlefield edge belonging to its deployment zone.

\subsubsection{Orders}

\begin{description}
    \item[Advance:] The Flyer performs a ground-attack mission. It may move and use its weapons normally. The detachment receives \textit{Protection (4+)}.

    \item[Charge:] A detachment with this order may perform an Aerial Interception if it has the \textit{Interceptor} ability, or directly charge a detachment at altitude. The detachment receives \textit{Protection (3+)}.

    \item[Forced March:] The detachment performs evasive manoeuvres but may use its weapons with Snap Fire. This order also allows the Flyer to land. The detachment receives \textit{Protection (3+)}.
\end{description}

\subsubsection{Movement}

When activated during the Movement Phase, Flyers move in a straight line.

A Flyer at altitude has an unlimited Movement characteristic and may move any desired distance, subject to a minimum movement of 45~cm. It may only change its facing at the beginning of its movement.

While moving, a Flyer ignores all terrain features and zones of control and may end its movement over Impassable terrain.

Flyers are never pinned in an assault.

A Flyer's movement may take it off the battlefield. If this occurs, the detachment must remain off-table for one complete turn. It may then return by beginning its movement from any battlefield edge.

\subsubsection{Landing}

A Flyer with a Forced March order may land if it has a Movement characteristic or the \textit{Transport} ability.

To land, the Flyer moves at altitude until it is directly above the intended landing point and then descends to ground level. It cannot move any farther during that turn.

While landed, it is treated as a Heavy Antigrav and uses the Movement characteristic shown on its profile. It also loses its \textit{Dodge} ability while at ground level.

A Flyer may use its weapons with Snap Fire during the turn in which it lands.

During the turn after landing, it may choose one of the following actions:

\begin{itemize}
    \item Remain in place and act as a Heavy Antigrav for the turn.

    \item Take off with a Forced March order. The Flyer rises from its current position to altitude and then acts as a normal Flyer. Its movement begins from that point with any desired facing. It may land again elsewhere during the same movement if desired.
\end{itemize}

\subsubsection{Flyer Shooting}

Flyers may shoot normally during the Combat Phase if they have an Advance or Forced March order.

\subsubsection{Abort Mission}

A Flyer detachment that has not yet acted during the current turn may abort its mission to gain additional protection.

It may only do so if it has not yet:

\begin{itemize}
    \item Fired any of its weapons.
    \item Disembarked troops.
    \item Performed an Aerial Interception.
    \item Engaged in an assault.
    \item Landed during the current turn.
\end{itemize}

If the detachment aborts its mission, it cannot perform any further actions during the current turn but may complete its movement. All attacks targeting it suffer a $-1$ To-Hit penalty.

A mission may be aborted in reaction to a shooting attack, provided the decision is declared before any dice are rolled.

\subsubsection{Shooting at Flyers}

Flyers at altitude are difficult targets. A weapon must have a 360\textdegree{} firing arc to target them.

A Flyer's Protection save depends on its current order.

A Flyer that lands may be targeted by Overwatch Fire at the end of its movement as though it were a ground-level target.

\subsubsection{Bombing}

Flyers may use their bombs during movement, but only while at altitude.

\subsubsection{Transport}

A landed Flyer follows the normal Transport rules.

While a Flyer is at altitude, only bases with the \textit{Jump Packs} or \textit{Antigrav} ability may disembark. They may disembark at any point during the Flyer's movement.

Bases can never embark while the Flyer is at altitude.

If a transport Flyer is destroyed while at altitude, only transported bases with the \textit{Jump Packs} or \textit{Antigrav} ability may attempt an Emergency Exit Test. All other transported bases are automatically destroyed.

\subsubsection{Aerial Interception}

A Flyer detachment with a Charge order may enter Aerial Interception. When doing so, it must make a movement of exactly 45~cm.

It may then intercept an enemy detachment moving or arriving at altitude within 45~cm. This allows it to engage both enemy Flyers and other bases at altitude.

If the detachment performs an interception, pause the enemy detachment's movement. The intercepting Flyers move in a straight line to the interception point, and an aerial assault is resolved immediately.

The intercepting Flyers remain at the interception point. Any surviving intercepted bases then complete their movement.

If no enemy detachment was intercepted, the Flyer detachment may engage an enemy detachment at altitude within 45~cm at the end of the Movement Phase.

\subsubsection{Aerial Assault}

An aerial assault may be initiated through Aerial Interception or by directly engaging an enemy detachment at altitude.

Consolidation is resolved normally. Outnumbering dice are also generated normally if several aerial assaults are performed during the same turn through interceptions.

Flyers are never pinned and are not considered to remain engaged in an assault.

The aerial assault is resolved during the Combat Phase. If the target is neither a Flyer nor a Floater, it suffers a $-2$ AF penalty.

\subsubsection{Morale}

If a Flyer detachment fails a Morale Test, it immediately leaves the battlefield by moving in a straight line.

It must remain off-table for one complete turn and may then return from any battlefield edge.

